\subsection*{8. Digressions in two-dimensional geometry}
The eighth section, beginning on p.~45, gathers the more concrete
two-dimensional questions arising from teaching and research supervision.
On p.~46 Grothendieck describes a repeated sequence: a promising initial
exploration fails to become a shared undertaking with students, and
stronger interests then displace it. Plane polygons and their moduli
varieties form the first theme, including the unexpected appearance of
algebraic geometry. The next theme concerns immersed curves and surfaces
with boundary. The ``stable'' and ``cellulaire'' conditions lead back to
the maps of \S3, while the ``balayage'' of a branch across a double point
and the erasure or creation of a ``bigône'' supply the elementary
operations for a dynamic study (pp.~46--47). The stated questions concern
classes of immersions and the extensions of a prescribed boundary
immersion; the reported tangible result is a canonical dévissage into
indecomposable immersions by ``téléscopage''. His assessment changes with
the third theme on p.~47: work on systems of lines and pseudolines has
brought more productive exchanges with students, although their joint
work has remained elementary. On p.~48 he describes the recent
combinatorial construction of a ``surface modulaire'' parametrizing an
additional pseudoline's relative positions, or the system's
``affinisations'' through a choice of pseudoline at infinity. The
catalogue names the related archive dossier
\href{https://grothendieck.umontpellier.fr/archives-grothendieck/}{\emph{Système de pseudo-droites} (cote 154)}.
The final movement returns explicitly to \S3: since 1977, the conformal
realization of maps, their algebraic curves in the oriented compact case,
and the associated metric possibilities have remained behind these
combinatorial questions. Thurston's geometric ideas reinforce that
expectation. At the transition to p.~49, however, the account of his
students' difficulties leads to the judgment that teaching has often
burdened rather than stimulated his geometric reflection; this judgment
opens the broader appraisal in \S9.
(\href{https://grothendieck.umontpellier.fr/web/viewer.html?file=../119.pdf\#page=48}{Copy 1, pp.~45--49, beginning at site 47};
\href{https://grothendieck.umontpellier.fr/web/viewer.html?file=../119.pdf\#page=49}{elementary operations on pp.~46--47, beginning at site 48};
\href{https://grothendieck.umontpellier.fr/web/viewer.html?file=../119.pdf\#page=51}{moduli surface and conformal viewpoint on p.~48, site 50};
\href{https://grothendieck.umontpellier.fr/web/viewer.html?file=../119.pdf\#page=52}{transition to the teaching appraisal on p.~49, site 51}.)
